%
% Template for RBIE papers in LaTeX
%

% The above language combination is for this template document only.
% You should use one of the following:
%\documentclass[english, spanish, brazilian]{RBIEarticle} % for papers in portuguese
%\documentclass[brazilian, spanish, english]{RBIEarticle} % for papers in english
%\documentclass[brazilian, english, spanish]{RBIEarticle} % for papers in spanish

% Papers in Portuguese or Spanish may require the following lines:
%
% Template for RBIE papers in LaTeX
%

\documentclass[english, spanish, brazilian]{RBIEarticle}
\title{Repositórios de Código e Padronização Institucional: Um Estudo Automatizado a partir do GitHub}
\titleinenglish{Code Repositories and Institutional Standardization: An Automated Study Based on GitHub}
\titleinspanish{Repositorios de Código y Estandarización Institucional: Un Estudio Automatizado a partir de GitHub}

\author{%
	\parbox{4cm}{%
		Gilmar G Nascimento\\
		IFAM\\
		ORCID: \href{https://orcid.org/0009-0002-1319-9824}{0009-0002-1319-9824}\\
	%	\href{mailto:gilmar.nascimento@ifam.edu.br}{gilmar.nascimento@ifam.edu.br}
	}
	\parbox{4cm}{%
		Maria Claudia Emer\\
		UTFPR\\
		ORCID: \href{https://orcid.org/0000-0002-0963-1891}{0000-0002-0963-1891}\\
	%	\href{mailto:claudiaemer@utfpr.edu.br}{claudiaemer@utfpr.edu.br}
	}
	\parbox{4cm}{%
		Adolfo Gustavo SS Neto\\
		UTFPR\\
		ORCID: \href{https://orcid.org/0000-0002-0260-5922}{0000-0002-0260-5922}\\
	%	\href{mailto:adolfogustavo@utfpr.edu.br}{adolfogustavo@utfpr.edu.br}
	}
	\parbox{4cm}{%
		Laudelino C Bastos\\
		UTFPR\\
		ORCID: \href{https://orcid.org/0000-0002-1585-7168}{0000-0002-1585-7168}\\
	%	\href{mailto:laudelinobastos@utfpr.edu.br}{laudelinobastos@utfpr.edu.br}
	}
}

\heading{Nascimento, G. G., Emer, M. C., Seca Neto, A. G. S., \& Bastos, L. C.}{RBIE v.34 – 2026}

\citeas{Nascimento, G. G., Emer, M. C., Seca Neto, A. G. S., \& Bastos, L. C. (2026). Repositórios de Código e Padronização Institucional: Um Estudo Automatizado a partir do GitHub. Revista Brasileira de Informática na Educação, 34. \url{https://doi.org/10.5753/rbie.2026.xxxx}}

\Submission{20/05/2026}
\CameraReady{}
\FirstRoundNotif{}
\EditionReview{}
\NewVersion{}
\AvailableOnline{}
\SecondRoundNotif{}
\Published{}\begin{document}
	
	\maketitle
	
% If the manuscript is written in English, then this element must be removed.
\begin{otherlanguage}{brazilian}
\begin{abstract}
	Este artigo analisa a adoção institucional de plataformas de repositórios de código nos Institutos Federais (IFs) da Amazônia Legal. A partir de uma busca exploratória no GitHub e da fundamentação nos conceitos de Engenharia de Software e referatório, identifica-se uma lacuna de institucionalização de práticas sistemáticas de versionamento e documentação no ambiente educacional e de pesquisa da região. Correlacionam-se os padrões estabelecidos em Projetos Pedagógicos Institucionais (PPIs) e Projetos Pedagógicos de Curso (PPCs) com as práticas exigidas pela Engenharia de Software, buscando uma amálgama entre a produção de códigos — em ensino, pesquisa e extensão — e a infraestrutura institucional de repositórios. Adicionalmente, são discutidas observações sobre código aberto, ciência aberta, versionamento, conformidade com a LGPD e a relação entre Ambientes Virtuais de Aprendizagem (AVAs) e plataformas versionadoras de código. Os resultados indicam que a maioria dos IFs da região analisados não possui organização oficial no GitHub, e uma parcela significativa sequer está registrada na plataforma. Como contribuição, propõe-se o desenvolvimento de recomendações baseadas no conceito de referatório — um instrumento de curadoria educacional — para a organização e aproveitamento pedagógico e científico de repositórios de código no contexto da Amazônia Legal, com potencial de escala para toda a Rede Federal.
\keywords <Repositórios de código, Institutos Federais. Amazônia Legal. Referatório. Ciência Aberta
\end{abstract}

\end{otherlanguage}

\begin{otherlanguage}{english}
\begin{abstract}
<Here comes the abstract of the paper in English. The abstract should summarize the contents of the manuscript and should contain at least 150 and at most 300~words long and must be written in italics, Times~10, justified, with no special indentation and no spacing before or after.>
\keywords <Abstract must be followed by 3 to 10 keywords. The keywords should be justified with a line space single, no special indentation, with no spacing before and spacing of exactly 24-points after. The text should be set in Times 10-point font size and in italic font style. Please use semi-colon as a separator. Keywords must be title cased.>
\end{abstract}
\end{otherlanguage}

% If the manuscript is written in English, then this element must be removed.
\begin{otherlanguage}{spanish}
\begin{abstract}
<Aquí viene el resumen del artículo en español. El resumen debe resumir el contenido del manuscrito y debe contener un mínimo de 150 y un máximo de 300~palabras y debe estar escrito en cursiva, Times~10, justificado, sin sangría especial y sin espacio antes o después.>
\keywords <El resumen debe ir seguido de 3 a 10 palabras clave. Las palabras clave deben estar justificadas con un espacio de línea simple, sin sangría especial, sin espacios antes y con un espacio de exactamente 24 puntos después. El texto debe configurarse fuente Times con tamaño de 10 puntos y en estilo de fuente cursiva. Utilice punto y coma como separador. Las palabras clave deben comenzar con una letra mayúscula.>
\end{abstract}
\end{otherlanguage}

		
%	
\section{Introdução}
	
	É possível escrever códigos utilizando apenas editores de texto simples e ferramentas de linha de comando isoladas. No entanto, essa abordagem demanda tempo excessivo e dispersa funcionalidades que poderiam estar agrupadas. Ambientes Integrados de Desenvolvimento (IDEs) surgiram exatamente para resolver esse problema: integram editor, compilador, depurador e versionamento em uma única interface, aumentando a produtividade e reduzindo a complexidade cognitiva.
	
	Essa mesma lógica de integração — agrupar funcionalidades dispersas para aumentar produtividade e reduzir complexidade — pode ser observada no âmbito institucional. A Rede Federal de Educação Profissional, Científica e Tecnológica (Rede Federal de EPCT) não é uma coleção solta de escolas técnicas. Há mais de um século, ela integra institutos, campi, cursos e níveis de ensino em uma estrutura nacional coordenada \cite{lei11922_2008}, potencializando o desenvolvimento tecnológico do país \cite{pacheco2011}.
	
	A integração é reconhecida como benéfica tanto para ferramentas de codificação quanto para a organização da educação profissional. No entanto, os códigos produzidos no âmbito dos Institutos Federais — instituições cujo próprio nome evoca \textbf{Educação, Ciência e Tecnologia} — permanecem, em grande medida, dispersos e sem política institucional de curadoria. Propõe-se que, assim como Projetos Pedagógicos Institucionais (PPIs) e Projetos Pedagógicos de Curso (PPCs) integram documentalmente a estrutura curricular\footnote{Documentos obrigatórios exigidos pelo Ministério da Educação (MEC) que fundamentam as atividades de ensino, pesquisa e extensão no âmbito do Sinaes.}, os repositórios de código sejam tratados como parte dessa mesma lógica de integração sistêmica, independentemente da área do conhecimento ou da ferramenta utilizada.
	
	O desenvolvimento de software é um processo complexo e necessita de profissionais bem preparados. Como argumenta Humphrey \cite{barr2018}, a sociedade contemporânea depende demasiadamente dos produtos de software para que possamos continuar com práticas artesanais em sua produção. A qualidade do software produzido está diretamente relacionada à formação recebida e às práticas adotadas durante o aprendizado \cite{sommerville2011}. A formação de profissionais — em qualquer área que envolva produção de código — deve incluir práticas eficazes de engenharia de software, aprendidas e exercitadas durante a educação formal. Nesse sentido, o domínio de ferramentas de versionamento e colaboração, como repositórios Git, torna-se competência fundamental não apenas para estudantes de computação, mas para todo aquele que produz código como parte de sua atividade acadêmica ou profissional.
	
	O termo ``repositório de software"~ abrange todos os artefatos produzidos durante o desenvolvimento de um sistema, desde arquivos com código-fonte armazenados em sistemas de controle de versão até mensagens trocadas entre desenvolvedores \cite{sokol2013}. Sobre o ensino e desenvolvimento de códigos na academia, Weinberg \cite{barr2018} observa que projetos de códigos feitos nas universidades geralmente não possuem manutenção, não são utilizados ou testados por outras pessoas. Parte significativa dos códigos derivados de práticas de ensino são ``engavetados", sem possibilidade de aproveitamento de boas iniciativas.
	
	Diferentemente do que ocorre com PPI e PPCs — cuja ausência inviabiliza o reconhecimento institucional perante o MEC — não há qualquer exigência normativa relacionada a repositórios de código. Um curso pode ser legalmente reconhecido sem que a instituição mantenha qualquer política de versionamento ou colaboração por meio de plataformas de hospedagem de código.
	
%	No entanto, a ausência de obrigação legal não equivale à ausência de relevância pedagógica e institucional. A Engenharia de Software, como disciplina, estabelece que o desenvolvimento sistemático de software envolve não apenas a codificação, mas um conjunto de atividades inter-relacionadas — especificação, projeto, implementação, testes, documentação, configuração e manutenção — que devem ser registradas e gerenciadas ao longo de todo o ciclo de vida \cite{sommerville2011,pressman2016}. Práticas como controle de versão, revisão de código por pares, rastreabilidade de requisitos, documentação arquitetural e automação de testes são componentes centrais de uma formação sólida \cite{pressman2016, fowler2018}. Essas práticas não se restringem à computação: um projeto de pesquisa em bioinformática, uma simulação numérica em engenharia ou um sistema de coleta de dados em ciências ambientais igualmente demandam versionamento, documentação e reprodutibilidade.
%	
%	Assim como é prática recomendada — ainda que nem sempre seguida — que um docente respeite a ementa de sua disciplina, a institucionalização de repositórios de código deveria ser uma prática esperada em toda a instituição que se nomeia como de ``Educação, Ciência e Tecnologia". Não se trata apenas de armazenar código, mas de \textbf{documentar o processo de construção e evolução} de artefatos de software, alinhando a formação e a pesquisa às práticas consolidadas da ciência aberta e da indústria \cite{sommerville2011}.
	
	Entre as plataformas que hospedam código aberto, destaca-se o GitHub. Segundo Borges \cite{borges2016}, seus repositórios concentram-se majoritariamente em bibliotecas e arcabouços para desenvolvimento web, ferramentas de código e componentes não web.\footnote{Embora o GitHub seja a plataforma mais utilizada e tenha sido o foco deste estudo por oferecer API pública para coleta automatizada de dados, existem alternativas como GitLab (auto-hospedado) e Codeberg (sem fins lucrativos). A análise comparativa dessas plataformas no contexto dos IFs, incluindo conformidade com a LGPD, é sugerida como trabalho futuro.} O potencial educacional e científico dessas plataformas, no entanto, depende de uma condição fundamental: a adoção institucional.
	
	Como será demonstrado adiante, a maioria dos Institutos Federais da Amazônia Legal analisados não possui organização oficial no GitHub, e uma expressiva parcela sequer está registrada na plataforma — o que sugere não uma violação legal, mas uma \textbf{lacuna de institucionalização de práticas sistemáticas de engenharia de software} no ambiente educacional e de pesquisa da região. Essa lacuna torna-se ainda mais significativa quando confrontada com os princípios da ciência aberta, que preconizam a disponibilização não apenas de publicações e dados, mas também de código-fonte como parte integrante do conhecimento científico.
	
	O simples acesso a plataformas como GitHub, contudo, não é suficiente para garantir seu aproveitamento pedagógico e institucional. É necessário um agrupamento intencional de informações, documentações e códigos que oriente educadores, pesquisadores e estudantes na navegação por esse ecossistema. É nesse ponto que o conceito de \textbf{referatório} — discutido a seguir — oferece uma contribuição arquitetural relevante.
	
	\textbf{Cumpre esclarecer: esta pesquisa não versa sobre o ensino de versionamento ou sobre as vantagens técnicas do Git. Seu objeto é a escala institucional — as políticas, diretrizes e normas (ou sua ausência) que deveriam orientar a gestão de repositórios de código nos IFs. Trata-se de um trabalho de natureza político-pedagógica, não de um minicurso sobre ferramentas.}
	
	Esta pesquisa parte do seguinte problema: diante da ausência de políticas institucionais estruturadas para a gestão de repositórios de código nos IFs da Amazônia Legal, em que medida um referatório educacional pode contribuir para a organização e curadoria desses repositórios como parte de uma política institucional de ciência aberta? Como objetivos específicos, busca-se: (i) identificar as plataformas de código atualmente utilizadas — ainda que informalmente — por docentes, pesquisadores e discentes da região; (ii) analisar quais fatores institucionais e pedagógicos favorecem ou restringem o uso dessas plataformas; (iii) com base nos dados levantados e no conceito de referatório, desenvolver recomendações aplicáveis à realidade da região amazônica.
	
	A escolha da Amazônia Legal como recorte não é arbitrária: suas particularidades geográficas, de infraestrutura e de concentração de docentes, discentes e projetos de pesquisa e extensão a tornam um lócus privilegiado para um estudo piloto. Ao submeter recomendações à revisão por pares, em vez de restringi-las a eventos institucionais internos, busca-se legitimidade científica e transparência. A hipótese de trabalho é que soluções adequadas à realidade amazônica possam, no futuro, ser estendidas a outras regiões da Rede Federal de EPCT, contribuindo para uma política nacional de curadoria de repositórios de código.
	
	\section{Referencial teórico}
	
	\subsection{Engenharia de Software e Ensino de Programação}
	
	A engenharia de software desempenha diversos papéis na construção de projetos que envolvem códigos. Conforme Pressman \cite{pressman2016}, a engenharia de software é uma disciplina relacionada com todos os aspectos da produção de software, desde os estágios iniciais de especificação até sua manutenção. Além dos processos técnicos, relaciona-se também com gerenciamento de projeto e desenvolvimento de ferramentas e métodos. Para o ensino de computação, isso implica que não basta ensinar sintaxe de linguagens; é necessário também ensinar práticas de organização, documentação e versionamento \cite{sommerville2011}.
	
	Há, no entanto, uma diferença entre o ambiente de desenvolvimento de habilidades relacionadas à engenharia de software no espaço acadêmico e no empresarial. Conforme Barr \cite{barr2018}, quem escreve bastante código na faculdade não necessariamente fará "bons" códigos nas empresas. Essa lacuna entre a formação acadêmica e as práticas profissionais tem sido objeto de estudo em diversas pesquisas na área de educação em computação \cite{qian2017}. Uma das hipóteses para essa lacuna é a ausência, nos currículos, de disciplinas que integrem teoria e prática por meio de ferramentas reais de desenvolvimento colaborativo \cite{arcuri2020}.
	
	Um exemplo particularmente relevante dessa lacuna é o Software de Pesquisa. Von Flach e colaboradores \cite{vonflach2023} o definem como o software desenvolvido durante o processo de pesquisa, incluindo código fonte, algoritmos, scripts, fluxos de trabalho computacionais e executáveis. Esse tipo de software é parte integrante do ecossistema de pesquisa moderno e é amplamente utilizado nas Ciências e Engenharias para gerar resultados que servem como evidência em publicações científicas.
	
	Devido à complexidade envolvida no desenvolvimento de software e ao conhecimento de domínio especializado necessário, os próprios pesquisadores desenvolvem o software de pesquisa ou estão intimamente envolvidos com sua criação \cite{carver2007}. No entanto, eles nem sempre estão familiarizados com boas práticas de desenvolvimento de software sustentável \cite{vonflach2023}.
	
	Na Engenharia de Software, duas linhas de pesquisa voltadas para sustentabilidade se destacam. A primeira é a Sustentabilidade de Software, cuja preocupação central é a capacidade do software de perdurar ao longo do tempo. A segunda é a Engenharia de Software para Sustentabilidade, também conhecida como SE4S \cite{venters2021}. Para o ensino de programação, isso implica que avaliar a qualidade de algoritmos exige ir além da verificação funcional. É necessário também considerar práticas de versionamento, documentação e organização que garantam a longevidade e a reprodutibilidade do código produzido.
	
	Dessa forma, os trabalhos desenvolvidos nas disciplinas, assim como os projetos de ensino, pesquisa e extensão, poderiam com os devidos cuidados ser inseridos em repositórios institucionais. Com curadoria e planejamento pedagógico, esses materiais poderiam ser reaproveitados em outros processos, como a identificação de erros de segurança, a análise de qualidade de código e o desenvolvimento de recursos pedagógicos.
	
	\subsection{Versionamento e Rastros de Aprendizagem}
	
	Além de servirem como repositórios de código, plataformas como GitHub registram automaticamente o histórico de desenvolvimento por meio de commits, mensagens de log e evolução temporal dos arquivos. Esses rastros digitais constituem uma fonte valiosa de dados para a análise do processo de aprendizagem em programação.
	
	Estudos recentes têm explorado o uso de logs de programação para compreender dificuldades de aprendizado e fornecer feedback automatizado aos estudantes \cite{qian2017, nascimento2026logs}. Especificamente, o uso de logs capturados durante a codificação permite identificar padrões de erro, estratégias de correção e evolução conceitual ao longo do tempo \cite{nascimento2026logs}.
	
	Nesta pesquisa, embora o foco principal seja a análise da adoção institucional de repositórios, reconhece-se o potencial dos logs de versionamento como ferramenta complementar para investigações futuras sobre a eficácia do uso do GitHub no ensino de programação, conforme apontado em trabalhos anteriores \cite{nascimento2026logs}. Além disso, os logs são salvos em repositórios, o que proporciona não apenas a continuidade de pesquisas futuras, mas também o alinhamento com os princípios da Ciência Aberta: dados de aprendizado tornam-se rastreáveis, compartilháveis e passíveis de reuso por outros pesquisadores, desde que respeitadas as questões éticas e de privacidade.
	
	\subsection{Ciência Aberta e Repositórios de Código}
	
	A Ciência Aberta tem ganhado destaque como um movimento que busca tornar a pesquisa mais transparente, reprodutível e acessível. No contexto do ensino de computação, plataformas como GitHub e GitLab são reconhecidas como infraestrutura essencial para a Ciência Aberta \cite{momcheva2025}. Os princípios FAIR (Findable, Accessible, Interoperable, Reusable) devem ser aplicados não apenas a dados, mas também a códigos produzidos no âmbito acadêmico.
	
	Pesquisas recentes mostram que o uso de repositórios Git no ensino promove:
	\begin{itemize}
		\item Participação e colaboração em projetos de código aberto; (quando o software/projeto é de terceiros(mesmo sendo código aberto))
		\item Versionamento e rastreabilidade de contribuições;
		\item Integração com serviços como Zenodo para geração de DOIs;
		\item Baixa barreira de entrada, dada a familiaridade crescente com Git \cite{vanGeest2024}.
	\end{itemize}
	
	\subsection{Plataformas, Soberania de Dados e Implicações Legais}
	
	No entanto, a adoção de plataformas de repositórios por instituições públicas brasileiras não é isenta de controvérsias. O GitHub, embora dominante, é uma plataforma proprietária adquirida pela Microsoft em 2018. Seus termos de serviço permitem, em determinadas circunstâncias, a utilização de dados públicos para treinamento de modelos de inteligência artificial e outras finalidades comerciais. Para uma instituição pública que produz conhecimento financiado com recursos da sociedade, essa dependência levanta questões sobre soberania de dados, propriedade intelectual e conformidade com a Lei Geral de Proteção de Dados (LGPD) \cite{lgpd2018}.
	
	Alternativas como GitLab e Codeberg oferecem modelos diferentes. O GitLab permite instalação auto-hospedada (GitLab Community Edition), mantendo repositórios em servidores próprios sob controle direto da instituição. O Codeberg, por sua vez, é uma plataforma sem fins lucrativos, comprometida exclusivamente com software livre, sem vínculo com grandes corporações.
	
	A escolha da plataforma, portanto, não é meramente técnica. Envolve decisões estratégicas sobre onde a instituição deposita seus ativos digitais — e sob quais condições legais e comerciais. Uma política institucional de repositórios não pode ignorar essas dimensões, sob pena de, ao buscar abertura, comprometer a soberania e a segurança jurídica da instituição.
	
	Nesta pesquisa, optou-se por analisar o GitHub por dois motivos: (i) sua dominância no cenário educacional e (ii) a disponibilidade de API pública para coleta automatizada de dados. A discussão sobre alternativas auto-hospedadas ou sem fins lucrativos, embora relevante, foge ao escopo deste estudo empírico, sendo sugerida como trabalho futuro.
	
	\subsection{Ambientes Virtuais de Aprendizagem e Repositórios}
	
	A prática comum em instituições de ensino, incluindo os IFs, é a utilização de Ambientes Virtuais de Aprendizagem (AVAs) como Moodle e Google Classroom para recebimento e avaliação de códigos. O fluxo tradicional envolve o estudante desenvolvendo código localmente, compactando em arquivo .zip e enviando ao professor, que avalia manualmente.
	
	Este fluxo, embora funcional, apresenta limitações:
	\begin{itemize}
		\item Não estimula boas práticas de versionamento (Git);
		\item Não promove reuso de código entre turmas e períodos;
		\item Não gera visibilidade para o trabalho do estudante;
		\item Não se alinha aos princípios de Ciência Aberta e reprodutibilidade.
	\end{itemize}
	
	Pesquisas indicam que existem ferramentas de integração entre AVAs e GitHub/GitLab. O Moodle, por exemplo, conta com plugins como VPL (Virtual Programming Lab) e CodeGrade que permitem submissão via \texttt{git push} e avaliação automática \cite{moodle2024}. O Google Classroom também oferece APIs e add-ons que possibilitam integração com GitHub \cite{google2024}.
	
	
	\subsection{Referatório}	
	O conceito de \textbf{referatório} é recente na literatura brasileira de educação digital, mas tem ganhado relevância como alternativa à simples acumulação de recursos educacionais. Um referatório pode ser concebido como um instrumento didático-metodológico que propicia o acesso qualificado a fontes relevantes de informação, favorecendo o desenvolvimento da autonomia intelectual e da competência informacional \cite{cecead_referatorio}.
	
	Diferentemente de um repositório tradicional, que armazena os recursos propriamente ditos (arquivos, códigos, documentos), um referatório organiza o acesso a eles. Como define a UniRede, ``um Referatório é um site na web que não faz o armazenamento dos recursos propriamente ditos, mas organiza o acesso a repositórios que detêm recursos sobre determinado assunto" \cite{unirede_referatorio}. Essa distinção é crucial: o referatório atua como uma camada de curadoria sobre infraestruturas existentes.
	
	No contexto educacional, o referatório pode funcionar como suporte para projetos interdisciplinares, sequências didáticas ou roteiros de aprendizagem, assumindo papel estratégico na mediação entre o conhecimento sistematizado e os sujeitos em formação \cite{cecead_referatorio}. Em tempos de sobrecarga informacional e abundância de conteúdos dispersos na internet, o referatório representa um contraponto fundamentado e organizado, pautado por critérios de relevância, confiabilidade e pertinência pedagógica.
	
	Aplicado ao universo de repositórios de código, o conceito ganha contornos específicos. As plataformas de hospedagem de código — GitHub, GitLab, Codeberg — são, em si mesmas, repositórios massivos, mas sua organização é descentralizada e dependente de ações individuais. Um professor que deseja encontrar exemplos de código produzidos por outros docentes da própria instituição não tem, hoje, um ponto de entrada único e curado. O referatório preencheria essa lacuna: não armazenaria os códigos (eles permaneceriam nos repositórios originais), mas organizaria o acesso qualificado a eles, agregando metadados pedagógicos, descrições de contexto de uso, nível de dificuldade, área do conhecimento e outros filtros relevantes para educadores e pesquisadores.
	
	Projetos já implementados na Rede Federal demonstram a viabilidade do conceito. O IFMG — Campus Ouro Preto mantém um referatório de objetos de aprendizagem \cite{ifmg_referatorio}, e a UniRede coordena um referatório de objetos de aprendizagem da EaD pública brasileira \cite{unirede_referatorio}. Ambos os casos indicam que a arquitetura de referatório é factível e pode ser adaptada para o contexto de repositórios de código.
	
	Transpondo essa lógica para o universo dos Institutos Federais, um referatório de repositórios educacionais de código poderia:
	\begin{itemize}
		\item Agregar repositórios de código produzidos por docentes, discentes e pesquisadores dos IFs da Amazônia Legal;
		\item Oferecer filtros por área do conhecimento, nível de ensino, tipo de atividade (ensino, pesquisa, extensão), licença de uso e tags pedagógicas;
		\item Integrar-se a políticas institucionais de ciência aberta, exigindo como condição de inclusão a adoção de licenças livres e a disponibilização pública do código;
		\item Servir como vitrine de boas práticas e repositórios de referência para a comunidade acadêmica regional.
	\end{itemize}
	
	O referatório, portanto, não é uma solução técnica para o problema de versionamento — que já é resolvido por Git e plataformas correlatas — mas uma \textbf{solução político-pedagógica} para o problema de visibilidade, curadoria e institucionalização dos repositórios de código produzidos no âmbito dos IFs.
	
	\section{Metodologia}
	
	A pesquisa adota uma abordagem mista, combinando procedimentos quantitativos e qualitativos. Foram analisados os oito Institutos Federais localizados na Amazônia Legal: IFAC (Acre), IFAP (Amapá), IFAM (Amazonas), IFPA (Pará), IFRO (Rondônia), IFRR (Roraima), IFMT (Mato Grosso) e IFTO (Tocantins). O IFMA (Maranhão), embora parcialmente inserido na Amazônia Legal, foi excluído da análise por estar majoritariamente fora da região.
	
	\subsection{Busca em plataformas de repositórios}
	
	Realizou-se uma busca exploratória na plataforma GitHub, utilizando os identificadores oficiais de cada IF (ex: IFTO, IFAC, IFAP, IFAM, IFPA, IFRR, IFRO, IFMT). A string de busca foi:
	
	\texttt{(ifto OR ifac OR ifap OR ifam OR ifpa OR ifrr OR ifro OR ifmt)}
	
	Foram registradas as seguintes informações para cada IF:
	\begin{itemize}
		\item Existência de organização oficial (ex: \texttt{github.com/ifto});
		\item Existência de pelo menos um repositório público associado ao identificador;
		\item Atividade nos últimos 12 meses (commits, issues, pull requests);
		\item Descrição e finalidade declarada dos repositórios (ensino, pesquisa, extensão, administrativo).
	\end{itemize}
	
	Ressalta-se que a análise concentrou-se na presença ou ausência de organizações institucionais, não na qualidade dos repositórios ou no domínio técnico de ferramentas por parte de docentes ou discentes. O interesse recai sobre a dimensão político-institucional — políticas, diretrizes, normativas — não sobre práticas individuais de versionamento.
	
	\subsection{Análise documental}
	
	Foram examinados os documentos institucionais publicamente disponíveis nos sites oficiais de cada IF, incluindo:
	\begin{itemize}
		\item Projetos Pedagógicos Institucionais (PPIs);
		\item Projetos Pedagógicos de Curso (PPCs) de cursos de computação e áreas afins;
		\item Planos de Desenvolvimento Institucional (PDIs);
		\item Relatórios de gestão (2023-2025);
		\item Políticas institucionais de ciência aberta ou de gestão de dados de pesquisa.
	\end{itemize}
	
	A análise buscou identificar menções explícitas a: (i) políticas de repositórios de código; (ii) diretrizes para versionamento e documentação de software produzido institucionalmente; (iii) incentivos à publicação de código aberto; (iv) integração entre AVAs e plataformas de código.
	
	\subsection{Análise da literatura}
	
	Realizou-se uma busca na Biblioteca Digital da SBC (SOL) e nos anais dos eventos SBIE, EduComp, CBIE e SBES dos últimos cinco anos, utilizando combinações dos termos: ``repositório de código", ``GitHub educacional", ``ensino de programação", ``ciência aberta", ``referatório"~ e ``institutos federais". Foram considerados trabalhos que abordassem o uso institucional de plataformas de código no contexto da educação profissional e tecnológica.
	
	\section{Resultados e Discussão}
	
	\subsection{Presença institucional no GitHub}
	
	A busca exploratória revelou que, dos oito IFs analisados:
	\begin{itemize}
		\item Apenas dois possuem organização oficial no GitHub (IFTO e IFPA);
		\item Quatro IFs não possuem qualquer organização registrada com seu identificador oficial;
		\item Um IF sequer aparece na plataforma em qualquer modalidade (nem organização, nem repositórios associados a docentes/discentes com o identificador);
		\item Os repositórios encontrados, quando existentes, estão majoritariamente associados a contas pessoais de docentes ou discentes, não a organizações institucionais.
	\end{itemize}
	
	Esses números sugerem que, embora haja atividade na plataforma (os 183 repositórios identificados inicialmente na busca ampla), ela não é resultado de uma política institucional estruturada, mas sim de iniciativas individuais.
	
	\subsection{Análise documental}
	
	A análise dos documentos institucionais revelou que nenhum dos IFs da Amazônia Legal possui, em seus PPIs, PDIs ou relatórios de gestão, menção explícita a:
	\begin{itemize}
		\item Políticas de repositórios de código;
		\item Diretrizes para versionamento de software produzido institucionalmente;
		\item Incentivos à publicação de código aberto como parte da produção acadêmica;
		\item Integração entre repositórios de código e AVAs institucionais.
	\end{itemize}
	
	Alguns IFs mencionam, em seus PPCs de cursos de computação, o ensino de ferramentas de versionamento como conteúdo programático. No entanto, não há qualquer menção a políticas institucionais que vão além do ensino da ferramenta — ou seja, que tratem o código produzido na instituição como patrimônio a ser curado, preservado e compartilhado.
		
	\subsection{Análise da literatura}
	
	Para compreender o estado da arte sobre o uso de repositórios de código no ensino de computação, realizou-se uma busca na Biblioteca Digital da SBC (SOL) utilizando a string \texttt{``repositório" and ``ensino"}. A busca retornou os trabalhos apresentados na Tabela~\ref{tab:sbc_literatura}.

	\begin{table}[!htb]
		\centering
		\small
		\caption{Trabalhos encontrados na SOL/SBC com a string ``repositório" and ``ensino" (versão resumida)}
		\label{tab:sbc_literatura_resumo}
		\begin{tabular}{p{12cm} c c c}
			\toprule
			\textbf{Título} & \textbf{Ano} & \textbf{GitHub} & \textbf{LMS} \\
			\midrule
			Identificação de Pesquisas  e Análise de Algoritmos de Clusterização para a Descoberta de Perfis de Engajamento & 2022 & Não & Sim \\
			Sobre a necessidade de recursos educacionais para o ensino do Pensamento Computacional na Educação Básica Brasileira: discussão e concepção de Repositório Educacional do Pensamento Computacional. & 2021 & Não & Não \\
			Small Open Online Course  com Professores do Ensino Médio: desafios para integrar REA nos materiais e atividades didáticas & 2016 & Não & Não \\
			Impacto da descontinuidade da  tecnologia Flash na disponibilidade dos Objetos de Aprendizagem em  repositórios educacionais & 2021 & Não & Não \\
			Versionamento de Projeto na Prática com Git e GitHub: Um Relato de Experiência do Curso Ofertado pelo Projeto LearningLab No Interior Cearense. & 2024 & Sim & Não \\
			Jogos Digitais sob a perspectiva da Engenharia de Software: um mapeamento sistemático da literatura & 2024 & Não & Não \\
			Implementação das Leis 10.639/03 e 11.645/08 no Ensino de Computação: Um Mapeamento da Literatura em Bases Nacionais & 2025 & Não & Não \\
			OntoObADi: Uma Ontologia para Repositório Web Semântico de Objetos de Aprendizagem Digitais & 2025 & Não & Não \\
			Uma API para Gerenciamento de Objetos de Aprendizagem Digitais Baseada em Web Semântica & 2025 & Não & Não \\
			Commit Explorer: Uma Solução Open Source para Extração e Avaliação de Commits e Códigos
			 & 2025 & Sim & Não \\
			Uma Análise Comparativa entre Repositórios de Recursos Educacionais Abertos para a Educação Básica & 2021 & Não & Não \\
			\bottomrule
		\end{tabular}
	\end{table}
	
	Dos treze trabalhos identificados, apenas dois mencionam diretamente o uso de GitHub em contexto educacional. O primeiro, "Versionamento de Projeto na Prática com Git e GitHub" (2024), é um relato de experiência sobre um curso de versionamento, ou seja, ensina a ferramenta em si. O segundo, "Commit Explorer" (2025), apresenta uma ferramenta para extração e avaliação de commits, com foco técnico. Nenhum dos dois, no entanto, aborda a dimensão institucional que esta pesquisa investiga: políticas de curadoria, integração com AVAs ou repositórios institucionais de código.
	
	Os demais trabalhos tratam de objetos de aprendizagem, repositórios educacionais para pensamento computacional, ou mapeamentos sistemáticos sem relação com versionamento de código. Observa-se também que nenhum dos trabalhos analisa a integração entre repositórios de código e Ambientes Virtuais de Aprendizagem (LMS) como Moodle ou Google Classroom, lacuna que esta pesquisa busca endereçar.
	
	A ausência de trabalhos que articulem repositórios de código, ensino e política institucional reforça a relevância deste estudo. A maioria das publicações encontradas ou ensina a ferramenta Git como conteúdo programático, ou propõe repositórios de objetos de aprendizagem no sentido tradicional (armazenamento de arquivos PDF, vídeos, apresentações). Não foram identificados trabalhos que discutam o referatório como solução de curadoria para códigos produzidos no âmbito dos Institutos Federais.
	
	\subsection{Síntese dos resultados}
	
	Os resultados indicam três constatações principais:
	
	\begin{enumerate}
		\item \textbf{Lacuna normativa:} não há exigência legal para repositórios de código, ao contrário de PPIs e PPCs;
		\item \textbf{Lacuna institucional:} a maioria dos IFs da Amazônia Legal não possui organização oficial no GitHub, e os repositórios existentes são fruto de ações individuais;
		\item \textbf{Lacuna de curadoria:} mesmo nos casos em que há repositórios, não existe uma camada de organização pedagógica que permita sua reutilização por outros docentes, pesquisadores ou discentes.
	\end{enumerate}
	
	Essas três lacunas convergem para um mesmo ponto: a ausência de uma \textbf{política institucional de gestão de repositórios de código}. É nesse ponto que o referatório se apresenta como uma solução arquitetural viável.
	
	\section{Considerações Finais}
	
	Esta pesquisa identificou uma lacuna institucional significativa nos IFs da Amazônia Legal: a ausência de políticas estruturadas para a gestão de repositórios de código. Embora não haja exigência normativa para tais políticas (ao contrário do que ocorre com PPIs e PPCs), a relevância pedagógica e científica da curadoria de códigos é inegável, especialmente para instituições de ``Educação, Ciência e Tecnologia".
	
	A contribuição principal deste trabalho é a proposição do referatório como instrumento de curadoria educacional e científica. O referatório não resolve o problema técnico de versionamento (já resolvido por Git e plataformas correlatas), mas sim o problema político-institucional de visibilidade, organização e reutilização dos códigos produzidos na instituição. Trata-se, portanto, de uma contribuição no campo da gestão educacional, não da engenharia de software ou do ensino de ferramentas.
	
	Como recomendações aplicáveis à realidade amazônica, sugere-se:
	\begin{itemize}
		\item Que os IFs da região considerem a criação de organizações institucionais em plataformas de código (GitHub, GitLab ou Codeberg);
		\item Que essas organizações sejam regidas por políticas claras de licenciamento, em conformidade com os princípios da ciência aberta;
		\item Que seja desenvolvido um referatório institucional ou interinstitucional, agregando os repositórios de código produzidos na região com curadoria pedagógica e científica;
		\item Que esse referatório sirva como piloto para uma política nacional da Rede Federal de EPCT.
	\end{itemize}
	
	Como limitações do estudo, reconhece-se que: (i) a busca concentrou-se no GitHub, não abrangendo exaustivamente outras plataformas; (ii) a análise documental baseou-se exclusivamente em fontes públicas; (iii) não foram realizadas entrevistas com gestores ou docentes para aprofundar as causas da lacuna identificada. Trabalhos futuros podem incluir a validação das recomendações junto a gestores da Rede Federal, o desenvolvimento de um protótipo de referatório para a região amazônica e a replicação do estudo em outras regiões do país.
	
	A hipótese que orienta este trabalho é que soluções adequadas à realidade amazônica — com suas particularidades geográficas, infraestruturais e de concentração de docentes, discentes e projetos — podem ser estendidas a outras regiões da Rede Federal. O artigo submetido à revisão por pares é o primeiro passo nessa direção: conferir legitimidade científica a uma discussão que, tradicionalmente, ocorreria apenas em reuniões institucionais internas.
	
	%\section{Referências}
\bibliographystyle{apacite}
\bibliography{references}
	
\end{document}